\subsection*{MR1--MR3: a numerical regularity bound for the actual genus-three input}
Let $X=J(C)$ for a smooth nonhyperelliptic genus-three curve, $L=3\Theta$, and $i:X\hookrightarrow\mathbf P^{26}$ the complete embedding. Its degree is $L^3=162$. In the $d=12$ construction the input is
\[
 F=I_{Z/X}(\Theta),\qquad Z=C_1\sqcup\cdots\sqcup C_{13},
\]
with thirteen disjoint translated Abel--Jacobi curves. The same bounds apply to the reflected copies used for the second input. The exact supplied-source locators are the genus-three seed, lines~3743--3776, the disjointness hypotheses, lines~5124--5135, and equation \texttt{eq-F-1-and-F-2}, lines~5403--5410.

\begin{proposition}[MR1--MR3: input regularity and presentation]\label{MR}
With this normalization,
\[
 P_C(t)=9t-2,\qquad P_Z(t)=117t-26,\qquad
 P_F(t)=27t^3+27t^2-108t-12.
\]
The sheaf $i_*F$ is $445$-regular. It has a presentation
\begin{equation}\label{eq:inputpresentation}
 0\longrightarrow K\longrightarrow
 \cO_X(-445L)^{\oplus 2384568978}\longrightarrow F\longrightarrow0,
\end{equation}
where $K$ is locally free and $i_*K$ is $449$-regular.
\end{proposition}
\begin{proof}
The Abel--Jacobi class is $\Theta^2/2$, so each curve has $L$-degree nine and Euler characteristic $-2$. This proves the first two polynomials. The defining sequence and Riemann--Roch give
\[
 \chi(F(tL))=(3t+1)^3-13(9t+1),
\]
which is the third polynomial.

The Gotzmann expansion of $9t-2$ consists of nine linear binomial terms, summing to $9t-27$, followed by twenty-five constants. Its length is $34$, so Gotzmann regularity, in the form \citet[Theorem~2.7]{Dellaca14}, bounds each ambient curve ideal by $34$. The thirteen curves are disjoint: at every ambient stalk at most one of their ideals is nonunit. Hence their tensor product equals their intersection. The product regularity bound of \citet[Proposition~1.5 and Corollary~1.6]{Sidman01} gives
\[
 \operatorname{reg}I_{Z/\mathbf P^{26}}\le13\cdot34=442.
\]
For comparison, using only the complete union polynomial gives the weaker Gotzmann number $6760$, since the $117$ linear terms sum to $117t-6669$ and $6643$ constants remain.

The sheaf $B=i_*\cO_X(\Theta)$ is $3$-regular: for $j=1,2,3$ its defining groups are $H^j(X,\cO_X((10-3j)\Theta))$, which vanish by positivity on an abelian variety. Higher groups vanish by support dimension. Put $A=I_{Z/\mathbf P^{26}}$. It is locally free off the one-dimensional $Z$, so Sidman's Proposition~1.5 gives $\operatorname{reg}(A\otimes B)\le445$. In that proposition the higher Tor sheaves are supported on $Z$; Lemma~1.4 applies to their dimension at most two. There is an exact sequence
\[
 0\longrightarrow T\longrightarrow A\otimes B\longrightarrow i_*F\longrightarrow0
\]
with $T$ supported on $Z$. For every $j\ge1$, $H^{j+1}(T(k))=0$ for all $k$. The cohomology sequence therefore passes $445$-regularity to $i_*F$.

Regularity gives global generation and higher-cohomology vanishing for $F(445L)$. Its number of sections is $P_F(445)=2384568978$, proving the surjection in \eqref{eq:inputpresentation}. The disjoint smooth curves form a codimension-two local complete intersection. Their ideal has projective dimension at most one, so the kernel of this locally free surjection is locally free.

Finally $i_*\cO_X$ is $4$-regular. For $j\ge2$ the sequence \eqref{eq:inputpresentation}, twisted by $(449-j)L$, proves the required vanishing for $K$. For $j=1$, regularity gives surjectivity of
\[
 H^0(F(445L))\otimes\operatorname{Sym}^3H^0(\mathbf P^{26},\cO(1))
 \longrightarrow H^0(F(448L)).
\]
It factors through the section map in \eqref{eq:inputpresentation} at twist $448L$. This proves the remaining vanishing and hence $449$-regularity.
\end{proof}

The input sheaf is now quantitatively controlled. Markman's Proposition~9.2.1 uses a direct-image complex $K^1\to K^2$ on $X\times\widehat X$ and identifies the output as the cokernel of its dual. Its active proof, supplied source lines~5846--5902, gives no ranks or twists for those two bundles. The next section performs the direct-image, normalization and isogeny estimates explicitly. The input regularity proposition by itself is not an output regularity estimate.
